\documentclass[10pt,a4paper]{amsart}
\usepackage{amsmath,amssymb,amsthm,graphicx}
\usepackage{verbatim}
\newenvironment{reference}{\comment}{\endcomment}
%\newenvironment{reference}{}{}
\newenvironment{slogan}{\comment}{\endcomment}
\newenvironment{history}{\comment}{\endcomment}

% For commutative diagrams we use Xy-pic
\usepackage[all]{xy}

% We use 2cell for 2-commutative diagrams.
\xyoption{2cell}
% Theorem environments.
%
\theoremstyle{plain}
\newtheorem{theorem}[subsection]{Theorem}
\newtheorem{proposition}[subsection]{Proposition}
\newtheorem{lemma}[subsection]{Lemma}

\theoremstyle{definition}
\newtheorem{definition}[subsection]{Definition}
\newtheorem{example}[subsection]{Example}
\newtheorem{exercise}[subsection]{Exercise}
\newtheorem{situation}[subsection]{Situation}

\theoremstyle{remark}
\newtheorem{remark}[subsection]{Remark}
\newtheorem{remarks}[subsection]{Remarks}

\numberwithin{equation}{subsection}

% Macros
%
\def\lim{\mathop{\mathrm{lim}}\nolimits}
\def\colim{\mathop{\mathrm{colim}}\nolimits}
\def\Spec{\mathop{\mathrm{Spec}}}
\def\Hom{\mathop{\mathrm{Hom}}\nolimits}
\def\Ext{\mathop{\mathrm{Ext}}\nolimits}
\def\SheafHom{\mathop{\mathcal{H}\!\mathit{om}}\nolimits}
\def\SheafExt{\mathop{\mathcal{E}\!\mathit{xt}}\nolimits}
\def\Sch{\mathit{Sch}}
\def\Mor{\mathop{\mathrm{Mor}}\nolimits}
\def\Ob{\mathop{\mathrm{Ob}}\nolimits}
\def\Sh{\mathop{\mathit{Sh}}\nolimits}
\def\NL{\mathop{N\!L}\nolimits}
\def\CH{\mathop{\mathrm{CH}}\nolimits}
\def\proetale{{pro\text{-}\acute{e}tale}}
\def\etale{{\acute{e}tale}}
\def\QCoh{\mathit{QCoh}}
\def\Ker{\mathop{\mathrm{Ker}}}
\def\Im{\mathop{\mathrm{Im}}}
\def\Coker{\mathop{\mathrm{Coker}}}
\def\Coim{\mathop{\mathrm{Coim}}}

% Boxtimes
%
\DeclareMathSymbol{\boxtimes}{\mathbin}{AMSa}{"02}

%
% Macros for moduli stacks/spaces
%
\def\QCohstack{\mathcal{QC}\!\mathit{oh}}
\def\Cohstack{\mathcal{C}\!\mathit{oh}}
\def\Spacesstack{\mathcal{S}\!\mathit{paces}}
\def\Quotfunctor{\mathrm{Quot}}
\def\Hilbfunctor{\mathrm{Hilb}}
\def\Curvesstack{\mathcal{C}\!\mathit{urves}}
\def\Polarizedstack{\mathcal{P}\!\mathit{olarized}}
\def\Complexesstack{\mathcal{C}\!\mathit{omplexes}}
% \Pic is the operator that assigns to X its picard group, usage \Pic(X)
% \Picardstack_{X/B} denotes the Picard stack of X over B
% \Picardfunctor_{X/B} denotes the Picard functor of X over B
\def\Pic{\mathop{\mathrm{Pic}}\nolimits}
\def\Picardstack{\mathcal{P}\!\mathit{ic}}
\def\Picardfunctor{\mathrm{Pic}}
\def\Deformationcategory{\mathcal{D}\!\mathit{ef}}
\usepackage{cleveref}
\crefname{originallemma}{original lemma}{original lemmas}
\Crefname{originallemma}{Original lemma}{Original lemmas}
\makeatletter
\newlabel{algebra-lemma-finite-dimensional-algebra}{{00J6}{source}{}{sourcepointer-1}{}}
\newlabel{algebra-lemma-finite-dimensional-algebra@cref}{{[originallemma][0][]00J6}{[1][0][]source}{}{}{}}
\makeatother
\crefname{originalproposition}{original proposition}{original propositions}
\Crefname{originalproposition}{Original proposition}{Original propositions}
\makeatletter
\newlabel{algebra-proposition-dimension-zero-ring}{{00KJ}{source}{}{sourcepointer-2}{}}
\newlabel{algebra-proposition-dimension-zero-ring@cref}{{[originalproposition][0][]00KJ}{[1][0][]source}{}{}{}}
\makeatother
\title{A linear relation between trace maps forces two surjections to agree}
\author{The Stacks Project Authors}
\date{stacks-project:weil.tex:local-snapshot-20261012; GNU FDL 1.2 or later}
\begin{document}
\maketitle
\noindent\textit{Editorial note.} All delivered mathematical text is the author's own, copied from the pinned Stacks Project chapter file weil.tex; no statement or proof step was rewritten, reordered or removed.\par\medskip
\noindent\textit{Editorial note.} The counted proof is the author's complete argument: the reduction to a finite separable algebra, the base change to an algebraic closure, and the minimal idempotent computation with the integrality of traces.\par\medskip
\noindent\textit{Editorial note.} The label lemma-algebra-relations names the counted statement, and the proof environment immediately following it in the source is delivered verbatim as its proof.\par\medskip
\noindent\textit{Editorial note.} The counted statement is carried by the author's own lemma environment in the chapter source.\par\medskip
\noindent\textit{Editorial note.} Pre-open gate: python operations/agent-authoring/tools/precard\_gate.py --source weil.tex --label lemma-algebra-relations --start 3590 --end 3653 --local-id 28131 returned RESULT=OK (L28131 free, occupancy FREE, public interval clear). Pre-assembly tag self-check: canonical tag resolved and not present in reference/public-items.jsonl.\par\medskip
\noindent\textit{Editorial note.} The statement and proof are copied from the pinned chapter file; the cited inputs are kept as external pointers with their canonical Stacks tags.\par\medskip
\noindent\textit{Allowed standard prerequisites (editorial).}\begin{itemize}
\item Decomposition of a finite dimensional algebra over a field into local Artinian factors. Algebra, the structure results for finite dimensional algebras and zero dimensional rings. Writes the finite algebra as a product of local Artinian algebras and reduces to the product of the separable residue fields.
\item Traces commute with base change of the coefficient field. The behaviour of the trace pairing of a finite separable algebra under an extension of the base field. Allows replacing the base field by an algebraic closure, where the separable algebras become products of copies of the field.
\item Structure of finite separable algebras over an algebraically closed field. The identification of a finite separable algebra over an algebraically closed field with a product of copies of the field, with its minimal idempotents. Reduces the relation to a computation on minimal idempotents, where the traces are integers.
\end{itemize}\medskip
\ifdefined\hypertarget\hypertarget{sourcepointer-1}{}\fi
\noindent\textit{Source reference.} Original lemma 00J6: sources/stacks-chow/tags.txt, source lines 704--704. Prerequisite pointer invoked at proof line(s) 3608; it names a standard fact of a sibling chapter that the counted argument uses, and every step of the counted proof built on it is retained verbatim in the delivered body.\par\medskip
\ifdefined\hypertarget\hypertarget{sourcepointer-2}{}\fi
\noindent\textit{Source reference.} Original proposition 00KJ: sources/stacks-chow/tags.txt, source lines 753--753. Prerequisite pointer invoked at proof line(s) 3609; it names a standard fact of a sibling chapter that the counted argument uses, and every step of the counted proof built on it is retained verbatim in the delivered body.\par\medskip
\begin{lemma}
\label{lemma-algebra-relations}
Let $F$ be a field of characteristic $0$.
Let $F'$ and $F_i$, $i = 1, \ldots, r$
be finite separable $F$-algebras. Let $A$ be a finite $F$-algebra.
Let $\sigma, \sigma' : A \to F'$ and $\sigma_i : A \to F_i$
be $F$-algebra maps. Assume $\sigma$ and $\sigma'$ surjective.
If there is a relation
$$
\text{Tr}_{F'/F} \circ \sigma - \text{Tr}_{F'/F} \circ \sigma' =
n(\sum m_i \text{Tr}_{F_i/F} \circ \sigma_i)
$$
where $n > 1$ and $m_i$ are integers, then $\sigma = \sigma'$.
\end{lemma}

\begin{proof}
We may write $A = \prod A_j$ as a finite product of
local Artinian $F$-algebras $(A_j, \mathfrak m_j, \kappa_j)$, see
Algebra, Lemma \ref{algebra-lemma-finite-dimensional-algebra} and
Proposition \ref{algebra-proposition-dimension-zero-ring}.
Denote $A' = \prod \kappa_j$ where the product is over those $j$
such that $\kappa_j/k$ is separable. Then each of the maps
$\sigma, \sigma', \sigma_i$ factors over the map
$A \to A'$. After replacing $A$ by $A'$ we may assume
$A$ is a finite separable $F$-algebra.

\medskip\noindent
Choose an algebraic closure $\overline{F}$. Set
$\overline{A} = A \otimes_F \overline{F}$,
$\overline{F}' = F' \otimes_F \overline{F}$, and
$\overline{F}_i = F_i \otimes_F \overline{F}$.
We can base change $\sigma$, $\sigma'$, $\sigma_i$
to get $\overline{F}$ algebra maps $\overline{A} \to \overline{F}'$
and $\overline{A} \to \overline{F}_i$. Moreover
$\text{Tr}_{\overline{F}'/\overline{F}}$ is the base
change of $\text{Tr}_{F'/F}$ and similarly for
$\text{Tr}_{F_i/F}$. Thus we may replace
$F$ by $\overline{F}$ and we reduce to the case discussed in
the next paragraph.

\medskip\noindent
Assume $F$ is algebraically closed and $A$ a finite separable $F$-algebra.
Then each of $A$, $F'$, $F_i$ is a product of copies of $F$.
Let us say an element $e$ of a product
$F \times \ldots \times F$ of copies of $F$ is a minimal idempotent
if it generates one of the factors, i.e., if
$e = (0, \ldots, 0, 1, 0, \ldots, 0)$. Let $e \in A$ be a minimal idempotent.
Since $\sigma$ and $\sigma'$ 
are surjective, we see that $\sigma(e)$ and $\sigma'(e)$ are minimal
idempotents or zero. If $\sigma \not = \sigma'$, then we can choose
a minimal idempotent $e \in A$ such that $\sigma(e) = 0$ and
$\sigma'(e) \not = 0$ or vice versa. Then
$\text{Tr}_{F'/F}(\sigma(e)) = 0$ and
$\text{Tr}_{F'/F}(\sigma'(e)) = 1$ or vice versa.
On the other hand, $\sigma_i(e)$ is an idempotent
and hence $\text{Tr}_{F_i/F}(\sigma_i(e)) = r_i$ is an integer.
We conclude that
$$
-1 = \sum n m_i r_i = n (\sum m_i r_i)
\quad\text{or}\quad
1 = \sum n m_i r_i = n (\sum m_i r_i)
$$
which is impossible.
\end{proof}

\end{document}
